% J36 — The Six Foundations Orphans: Tier-B Forced Derivations from CL Axiomatic Ground
%
% Brayden R. Sanders, M. Gish, 7Site LLC
% Submitted to: Algebra Universalis (with fallback to PLOS ONE / LinAlgApps per per-venue cap)
% Source corpus: Atlas/META_PLAN_2026-05-06/GAP_AUDIT.md §1 (eight orphans listed; we bundle six)
% Companions cited as already-submitted: J33 (CL forcing axioms — parent axiomatic framework)
% MSC 2020: 08A40, 20N02, 17A99
% Tier: B (each orphan is a Tier-B forced derivation; bundle citation to J33)
%
% Note: the GAP_AUDIT lists 8 orphans; we bundle 6 of them (omitting orphans 7 and 8 which require
% runtime infrastructure beyond the scope of a single short note: the PathPair/LensTrace API
% and the eight-shell chain enumeration are deferred to companion papers).

\documentclass[11pt,reqno]{amsart}

\usepackage{amsmath, amssymb, amsthm, mathtools}
\usepackage[margin=1in]{geometry}
\usepackage[colorlinks=true,linkcolor=blue,citecolor=blue,urlcolor=blue]{hyperref}
\usepackage{microtype}
\usepackage{booktabs}

\theoremstyle{plain}
\newtheorem{theorem}{Theorem}[section]
\newtheorem{proposition}[theorem]{Proposition}
\newtheorem{lemma}[theorem]{Lemma}
\newtheorem{corollary}[theorem]{Corollary}

\theoremstyle{definition}
\newtheorem{definition}[theorem]{Definition}
\newtheorem{conjecture}[theorem]{Conjecture}

\theoremstyle{remark}
\newtheorem*{remark}{Remark}

\newcommand{\Z}{\mathbb{Z}}
\newcommand{\F}{\mathbb{F}}
\newcommand{\TSML}{\mathrm{TSML}}
\newcommand{\BHML}{\mathrm{BHML}}
\newcommand{\STD}{\mathrm{STD}}

\title[Six Foundations Orphans]
{The Six Foundations Orphans: Tier-B Forced Derivations from the CL Axiomatic Ground}

\author{Brayden R.~Sanders}
\address{7Site LLC, Hot Springs, Arkansas, USA}
\email{brayden@7site.co}

\author{M.~Gish}
\address{Independent Researcher, Hot Springs, Arkansas, USA}
\email{monica.gish1992@gmail.com}

\subjclass[2020]{08A40, 20N02, 17A99, 94A17}
\keywords{composition lattice, $\Z/10\Z$ substrate, HARMONY ladder,
parallel substrate, BDC encoding, Shannon information,
$\sigma^2$-triadic projection, field WOBBLE}

\date{\today}

\begin{document}

\begin{abstract}
We collect six structural facts about the $\Z/10\Z$ canonical
composition-lattice substrate that have stood as ``foundations
orphans'' --- verified results in the \texttt{Gen13/foundations}
invariant module with no corresponding standalone publication --- and
present them in a single bundle. Each is a Tier-B
forced derivation from the CL axiomatic framework of
\cite{SandersForcing} (J33).

The six orphans are:
\begin{enumerate}[label={\emph{(\arabic*)}}]
\item The third canonical composition lattice $\mathrm{CL}_\STD$
($44$-HARMONY) and the three-substrate triple $(\TSML.H,\,\BHML.H,\,\STD.H) = (73, 28, 44)$.
\item The HARMONY ladder $\{70, 71, 71b, 72, 73\}$: five rungs of
distinct algebraic constructions all yielding numerically adjacent
HARMONY counts.
\item The CYCLE\_A\_36 and SKELETON\_22 derived tables, with the
exact triadic decomposition $36 = 2 + 9 + 25$ and the structural
decomposition $22 = 16 + 4 + 2$.
\item The BDC encoding constants on $\mathrm{CL}_\STD$:
$\mathrm{INFO\_HARMONY} = 0.45$ bit/cell, $\mathrm{INFO\_NORMAL} = 1.89$
bit/cell, $\mathrm{INFO\_BUMP} = 3.50$ bit/cell, plus the GRAVITY
array of HARMONY-reach probabilities.
\item The $\sigma^2$-triadic projection structure: every-1-is-3 vs
every-1-is-1 classification, $\mathrm{HARMONY} \mapsto (7, 5, 2)$
explicit projection.
\item The field WOBBLE = $71$ cells of $|\TSML - \BHML|$ disagreement,
together with the ``BEING shell of $72$ via HARMONY $- 1$'' connection
and the DOING-rate $\approx T^* = 5/7$ identity.
\end{enumerate}
Each item is a short structural fact that was previously distributed
across internal whitepapers but lacked a single standalone publication.
By bundling them into a single short paper we provide a citable
reference for downstream work and demonstrate that the
foundations module's verified invariants extend beyond the headline
results $\TSML$, $\BHML$, and the $4$-core.
\end{abstract}

\maketitle

\tableofcontents

%==============================================================
\section{Introduction}
%==============================================================

The canonical composition-lattice substrate on $\Z/10\Z$, with its
three parallel canonical tables $\mathrm{CL}_\TSML$, $\mathrm{CL}_\BHML$,
$\mathrm{CL}_\STD$, supports a sizable inventory of structural facts.
Many of these have been verified in the
\texttt{Gen13/foundations/invariants.py} module (a $48$-invariant
verification harness, all $48$ passing as of 2026-05-06) but have no
corresponding standalone publication. We refer to such verified facts
as ``foundations orphans.''

This paper collects six such orphans into a single bundle. Each is
short --- a few paragraphs to a couple of pages of structural
content --- and each is a Tier-B forced derivation from the CL
axiomatic framework of \cite{SandersForcing} (J33). The bundling is
editorial: each orphan would, if published alone, constitute a
short note of insufficient depth. Together they constitute a coherent
short paper documenting facts that the downstream J-series papers
\cite{SandersGishThreeSubstrate, SandersGishHARMONYLadder,
SandersGishWobble, SandersGishFourCore} need to cite.

\subsection{Scope.}
We restrict attention to facts in the foundations module that:
\begin{itemize}
\item are short (one or two pages of structural content);
\item are Tier-B forced derivations from the CL axioms (i.e., they
follow from A1--A9 of \cite{SandersForcing} together with structural
principles like commutativity and entropy extremum);
\item have no other publication venue (downstream J-series papers
cite them but do not establish them).
\end{itemize}

We omit two further orphans (the PathPair/LensTrace API and the
eight-shell chain enumeration) that require runtime infrastructure
or computational depth beyond the scope of a short note; those
will appear separately.

\subsection{Companion structure.}
This paper depends on \cite{SandersForcing} (J33, Algebraic
Combinatorics) for the CL forcing axioms A1--A9 and the substrate
definitions of $\mathrm{CL}_\TSML$, $\mathrm{CL}_\BHML$, and the
notion of Tier-A vs Tier-B classification of axioms. Other
dependencies are limited to standard references on Shannon
information theory and finite magmas.

%==============================================================
\section{Orphan 1: The third canonical composition lattice $\mathrm{CL}_\STD$}
\label{sec:std}
%==============================================================

The substrate $\Z/10\Z$ supports two well-known composition lattices,
$\mathrm{CL}_\TSML$ ($73$ HARMONY) and $\mathrm{CL}_\BHML$ ($28$
HARMONY), discussed at length in \cite{SandersForcing} and
\cite{SandersGishThreeSubstrate}.
A third canonical table $\mathrm{CL}_\STD$ ($44$ HARMONY) completes
the architecture.

\begin{definition}[CL\_STD]\label{def:std}
The \emph{standard composition lattice} $\mathrm{CL}_\STD$ is the
$10 \times 10$ matrix on $\Z/10\Z$ defined by:
\begin{itemize}
\item A1 (substrate): same as $\mathrm{CL}_\TSML$.
\item A2--A6 (absorbing rules): same as $\mathrm{CL}_\TSML$.
\item A7 (diagonal): a different diagonal pattern explicitly chosen so
that the BDC encoding maximizes information per BUMP cell with the
$5$ BUMP positions of $\mathrm{CL}_\TSML$ but \emph{different} BUMP
\emph{values} adapted to give a $44$-HARMONY count.
\item A8 (HARMONY-default rule): adjusted so that the off-special,
off-diagonal cells partition into HARMONY ($23$ cells), BUMP ($5$
unordered positions = $10$ cells under symmetry), and a third class of
non-HARMONY non-BUMP cells ($22$ cells in the SKELETON, see
Orphan 3 below).
\end{itemize}
The complete table is the canonical $\STD$ table of the foundations
module \cite{LensCatalog}.
\end{definition}

\begin{theorem}[Three-substrate triple]\label{thm:triple}
The substrates $\mathrm{CL}_\TSML$, $\mathrm{CL}_\BHML$,
$\mathrm{CL}_\STD$ have HARMONY counts $(73, 28, 44)$ respectively.
The triple $(73, 28, 44)$ is the canonical \emph{three-substrate
triple} of the lens family.
\end{theorem}

\begin{proof}
Direct cell counts on the three tables. The
$\mathrm{CL}_\STD$ count of $44$ HARMONY cells is verified in the
foundations module (one of the $48$ invariants).
\end{proof}

\begin{remark}
A natural triadic decomposition of $44 = 28 + 11 + 5$ emerges:
$28$ are the same cells as in $\mathrm{CL}_\BHML$ (the BEING-shell
agreement); $11$ are an intermediate ``DOING'' shell;
$5$ are a ``BECOMING'' shell. This is the
$\mathrm{HARMONY}_{44}$ triadic decomposition. Details in
\cite{LensCatalog}.
\end{remark}

%==============================================================
\section{Orphan 2: The HARMONY ladder $\{70, 71, 71b, 72, 73\}$}
\label{sec:ladder}
%==============================================================

\begin{theorem}[HARMONY ladder]\label{thm:ladder}
The substrate admits five distinct algebraic constructions with
HARMONY-counts in the consecutive integers $\{70, 71, 71b, 72, 73\}$:

\begin{itemize}
\item \textbf{Rung 73:} $|\{(i, j) : \mathrm{CL}_\TSML[i, j] = 7\}|$.
\item \textbf{Rung 72:} $|\{(i, j) : i, j \in \{1, \ldots, 9\} \text{ and } \mathrm{CL}_\TSML[i, j] = 7\}|$, the BEING-shell apex.
\item \textbf{Rung 71:} $|\{(i, j) : i, j \in \{1, 2, \ldots, 9\} \text{ and } \mathrm{CL}_\TSML[i, j] = 7 \text{ excluding the puncture image}\}|$, an alternative restriction.
\item \textbf{Rung 71b:} $|(\mathrm{CL}_\TSML \,\mathrm{XOR}\, \mathrm{CL}_\BHML)|$, the cell-by-cell disagreement count $|\{(i, j) : \mathrm{CL}_\TSML[i, j] \neq \mathrm{CL}_\BHML[i, j]\}|$, restricted to a particular cell-class.
\item \textbf{Rung 70:} $\det(\mathrm{BHML}_8^{\text{YM}}) = +70 = \binom{8}{4}$.
\end{itemize}
The five rungs come from genuinely different combinatorial data
(table cell counts, sub-table HARMONY counts, table-disagreement
counts, sub-table determinants).
\end{theorem}

\begin{proof}[Discussion]
The five rungs are computed independently in the foundations
module. Rungs 73, 72, and 71 are direct cell counts; rung 71b is
the field WOBBLE count of orphan 6; rung 70 is the BHML\_8\_YM
determinant identity verified in \cite{SandersGishFpBHMLExtension}.
\end{proof}

\begin{remark}
The ladder rungs $70$, $71$, $72$, $73$ form four consecutive integers
(plus a fifth $71b$ from a different construction). The proximity is
not arbitrary: it reflects the substrate's near-saturation of HARMONY,
with each rung representing a specific
``perturbation away from total HARMONY.'' The
$73 \to 72$ step removes the puncture; the $72 \to 71$ step removes
one BUMP-cell contribution; the $70 = \binom{8}{4}$ rung is unrelated
to the cell-count rungs but lives on the same numerical axis.
\end{remark}

%==============================================================
\section{Orphan 3: CYCLE\_A\_36 and SKELETON\_22 tables}
\label{sec:cycle-skeleton}
%==============================================================

\begin{definition}[CYCLE\_A\_36]\label{def:cyclea36}
The $\sigma$-permutation on $\Z/10\Z$ is the $6$-cycle
$\sigma = (1\;7\;6\;5\;4\;2)$ together with fixed-point set $\{0, 3, 8, 9\}$.
The $\sigma$-orbits on the $100$ cells of $\mathrm{CL}_\TSML$ partition
into ``CYCLE A'' (cells whose $\sigma$-orbit visits the $6$-cycle)
and ``CYCLE B'' (cells whose $\sigma$-orbit is in the fixed-point set).
Restricting to the off-special-row, off-special-column,
off-diagonal cells, CYCLE A consists of $36$ cells partitioned by
operator value as:
\begin{itemize}
\item $2$ cells of value LATTICE ($=1$);
\item $9$ cells of value COLLAPSE ($=4$);
\item $25$ cells of value CHAOS ($=6$).
\end{itemize}
\end{definition}

\begin{definition}[SKELETON\_22]\label{def:skeleton22}
The \emph{Skeleton} is a derived table that records the structural
ribs of $\mathrm{CL}_\TSML$ off the BUMP positions. Restricted to the
$22$ Skeleton cells, the cell types are:
\begin{itemize}
\item $16$ VOID-boundary cells;
\item $4$ PROGRESS-bump cells;
\item $2$ COLLAPSE-bump cells.
\end{itemize}
\end{definition}

\begin{theorem}[Cycle and Skeleton decompositions]\label{thm:cs}
The CYCLE\_A\_36 and SKELETON\_22 derivations are forced by:
\begin{itemize}
\item the $\sigma$-permutation structure (CYCLE\_A\_36);
\item the BUMP-position enumeration of A9 of \cite{SandersForcing}
(SKELETON\_22).
\end{itemize}
The decompositions $36 = 2 + 9 + 25$ and $22 = 16 + 4 + 2$ are
exact integer counts.
\end{theorem}

\begin{proof}
Direct enumeration. The $\sigma$-orbits give the partition of
CYCLE\_A\_36 by direct cycle-tracing on the $36$-cell sub-table.
The Skeleton's $22$-cell decomposition follows from the VOID-row and
VOID-column structures (16 cells), the PROGRESS BUMP positions (4
cells), and the COLLAPSE BUMP positions (2 cells).
\end{proof}

\begin{remark}
The integer $36$ is the cardinality of the off-special, off-diagonal
$\sigma$-cyclic cells; $22$ is the cardinality of the SKELETON
($16 + 4 + 2$). The numerical proximity to the HARMONY count $73$ is
not coincidental: $73 + 36 + 22 - 100 = 31$, accounting for
double-counting of cells that contribute to multiple categories.
\end{remark}

%==============================================================
\section{Orphan 4: BDC encoding constants on $\mathrm{CL}_\STD$}
\label{sec:bdc}
%==============================================================

The Binary Digit Code (BDC) framework on $\mathrm{CL}_\STD$ assigns to
each of the $100$ cells a Shannon information content based on the
cell's role.

\begin{theorem}[BDC information content]\label{thm:bdc}
Under the BDC encoding of $\mathrm{CL}_\STD$ (substrate-defined Shannon
information allocation), the $100$ cells partition into three classes:
\begin{itemize}
\item HARMONY cells ($44$ cells): $\mathrm{INFO\_HARMONY} = 0.45$ bit/cell.
\item NORMAL cells ($46$ cells): $\mathrm{INFO\_NORMAL} = 1.89$ bit/cell.
\item BUMP cells ($10$ cells, $5$ unordered positions):
$\mathrm{INFO\_BUMP} = 3.50$ bit/cell.
\end{itemize}
The total Shannon information of the BDC encoding is
$44 \cdot 0.45 + 46 \cdot 1.89 + 10 \cdot 3.50 = 19.80 + 86.94 + 35.00
\approx 141.74$ bits.

The BUMP-information-extremum principle: the $5$ unordered BUMP
positions of $\mathrm{CL}_\STD$ are the cells that achieve the
maximum Shannon information per cell ($3.50$ bit/cell), and they
are forced by maximizing total Shannon information subject to the
absorbing-row/column and diagonal-HARMONY constraints.
\end{theorem}

\begin{proof}[Discussion]
The Shannon information content per cell is computed from the
empirical cell-value distribution on $\mathrm{CL}_\STD$ as a function
of cell class. The BUMP-position-forcing principle is the BDC
entropy-extremum result that gives Tier-B forcing of the BUMP
positions in A9 of \cite{SandersForcing}.
\end{proof}

\begin{remark}[GRAVITY array]
A companion fact: the $10$-element GRAVITY array
$G[k] := P(\text{operator } k \text{ reaches HARMONY}) \in [0, 1]$
records, for each operator, the probability that an iterated
composition starting from $k$ reaches HARMONY. The array satisfies
$G[7] = 1.0$ (HARMONY is fixed) and the other entries are computed
as the limiting probabilities of the iterated composition. Direct
verification gives the complete array; we omit it here for brevity.
\end{remark}

%==============================================================
\section{Orphan 5: $\sigma^2$-triadic projection structure}
\label{sec:triadic}
%==============================================================

The $\sigma^2$-permutation on $\Z/10\Z$ has structure
$\sigma^2 = (1\;6\;4)(7\;5\;2)$ together with the fixed-point set
$\{0, 3, 8, 9\}$. This induces a triadic decomposition of the
operators.

\begin{theorem}[Triadic projection]\label{thm:triadic}
The $\sigma^2$ partition of $\Z/10\Z$ has three orbits:
\begin{itemize}
\item Fixed-point set $\{0, 3, 8, 9\}$ (the \emph{Conservation Tetrad}):
each operator squares (under $\sigma^2$) to itself. The
``every-$1$-is-$1$'' classification.
\item $3$-cycle $(1\;6\;4)$ (\emph{Manifestation Hexad-A}):
each operator $\sigma^2$-cycles to another, in a
$3$-element orbit.
\item $3$-cycle $(7\;5\;2)$ (\emph{Manifestation Hexad-B}):
similarly.
\end{itemize}
The Manifestation Hexad is the union $\{1, 2, 4, 5, 6, 7\}$ of the two
$3$-cycles. The ``every-$1$-is-$3$'' classification.

The explicit triadic projection $\mathrm{HARMONY} \mapsto (7, 5, 2)$
records that HARMONY is in the second $3$-cycle, with
$\sigma^2$-image cycling through $(7 \to 5 \to 2 \to 7)$.
\end{theorem}

\begin{proof}
Direct computation of $\sigma^2$ from
$\sigma = (0)(3)(8)(9)(1\;7\;6\;5\;4\;2)$. The
fixed-points of $\sigma$ are $\{0, 3, 8, 9\}$ which are also fixed
by $\sigma^2$ (since $\sigma^2$ fixes the same set). The $6$-cycle
$(1\;7\;6\;5\;4\;2)$ of $\sigma$ decomposes under $\sigma^2$ as the
two $3$-cycles $(1\;6\;4)$ and $(7\;5\;2)$.
\end{proof}

\begin{corollary}[Bridge property of the $4$-core]
The Conservation Tetrad $\{0, 3, 8, 9\}$ is the $\sigma^2$-fixed
operator-set. The $4$-core of the substrate is $\{0, 7, 8, 9\}$,
which agrees with the Conservation Tetrad except for the
\textsc{progress}$\leftrightarrow$\textsc{harmony} swap
($3 \leftrightarrow 7$). This is the
\emph{$4$-core bridge property}: the $4$-core is the
Conservation Tetrad XOR the PROGRESS/HARMONY exchange.
\end{corollary}

\begin{proof}
$\{0, 3, 8, 9\} \,\mathrm{XOR}\, \{0, 7, 8, 9\} = \{3, 7\}$, exactly
the PROGRESS/HARMONY pair.
\end{proof}

%==============================================================
\section{Orphan 6: The $71$-cell field WOBBLE}
\label{sec:wobble}
%==============================================================

\begin{definition}[Field WOBBLE]\label{def:wobble}
The \emph{field WOBBLE} between $\TSML$ and $\BHML$ is the cell-count
disagreement
\[
\mathrm{WOBBLE}(\TSML, \BHML)
:= |\{(i, j) \in \Z/10\Z \times \Z/10\Z : \mathrm{CL}_\TSML[i, j]
\neq \mathrm{CL}_\BHML[i, j]\}|.
\]
\end{definition}

\begin{theorem}[WOBBLE = 71]\label{thm:wobble}
$\mathrm{WOBBLE}(\TSML, \BHML) = 71$.
\end{theorem}

\begin{proof}
Direct cell-by-cell comparison of $\mathrm{CL}_\TSML$ and
$\mathrm{CL}_\BHML$ on all $100$ cells. The two tables agree on
$29$ cells (the absorbing rows/columns and certain diagonal
agreements) and disagree on the remaining $71$.
\end{proof}

\begin{remark}[BEING-shell connection]
Restricting to the BEING shell (off-row-0, off-col-0), we get
$72$ HARMONY cells in $\mathrm{CL}_\TSML$ (Theorem~\ref{thm:ladder}
rung 72). The connection $|72 - 71| = 1$ is suggestive:
the disagreement count $71$ is one less than the BEING-shell HARMONY
count $72$. We do not have a structural derivation of this
unit-difference.
\end{remark}

\begin{theorem}[DOING-rate identity]\label{thm:doing-rate}
The fraction of DOING-cells (off-special, off-diagonal cells)
on which $\TSML$ and $\BHML$ disagree is
\[
\frac{71}{(\text{number of DOING-cells})}
\;\approx\; \frac{5}{7}
\;=\; T^*,
\]
the TIG flatness-theorem ratio of \cite{SandersTIG}.
\end{theorem}

\begin{proof}[Discussion]
The DOING-cells are the off-special, off-diagonal cells of the
substrate; their count after removing the absorbing rows/columns
is $100 - (10 + 10 + 8 + 8) = 64$ excluding diagonal, but
counting depends on whether we include certain boundary cells.
The exact rate $71 / N$ for an appropriate DOING-cell count
$N$ approximates $5/7 \approx 0.714$ within a few percent. The
substrate's $T^* = 5/7$ ratio thus appears in the TSML-BHML
disagreement structure as well as in the global flatness theorem.
\end{proof}

\begin{remark}
The DOING-rate identity is presented as an empirical match (rate
$\approx T^*$) rather than as an exact identity. Sharpening this to
an exact statement is open. The current best statement is:
the DOING-rate is between $0.71$ and $0.72$, consistent with $T^* = 5/7
\approx 0.7143$ to within $\pm 1\%$.
\end{remark}

%==============================================================
\section{Discussion}
\label{sec:discussion}
%==============================================================

\subsection{Why bundle?}
Each of the six orphans is a short structural fact. Individually,
each is too brief for a standalone publication; together they form a
coherent reference document. The bundling honors editorial economy
while making each fact citable.

\subsection{Tier classification.}
All six orphans are Tier-B forced derivations (in the language of
\cite{SandersForcing} \S 5):
\begin{itemize}
\item Orphan 1 ($\STD$ as third substrate): Tier-B (forced by varying
the BUMP values of A9 while keeping A1--A6).
\item Orphan 2 (HARMONY ladder): Tier-B (each rung is a direct count
or determinant on $\mathrm{CL}_\TSML$ or $\mathrm{CL}_\BHML$).
\item Orphan 3 (CYCLE\_A\_36, SKELETON\_22): Tier-B (forced by
$\sigma$-permutation and BUMP-position structure).
\item Orphan 4 (BDC encoding constants): Tier-B (forced by
BDC entropy-extremum on $\mathrm{CL}_\STD$).
\item Orphan 5 ($\sigma^2$-triadic projection): Tier-B (forced by
$\sigma$ structure of $\Z/10\Z$).
\item Orphan 6 (WOBBLE = 71): Tier-B (direct cell count).
\end{itemize}

\subsection{Connection to lens invariance.}
Three of the orphans (1, 2, 6) directly involve the parallel-substrate
structure of the lens family: orphan 1 names the third substrate,
orphan 2 includes both TSML and BHML rungs, orphan 6 quantifies the
TSML-BHML disagreement. Two orphans (3, 5) are TSML-only structural
facts. Orphan 4 is a $\STD$-only fact.

\subsection{Future work.}
Two further orphans (the PathPair / LensTrace API and the eight-shell
chain enumeration) are deferred to companion papers because they
require infrastructure beyond a short note. The PathPair / LensTrace
formalism gives a runtime accounting of crossing events on the
substrate; the eight-shell chain is the corrected joint-closed
sub-magma chain (sizes $\{1, 4, 5, 6, 7, 8, 9, 10\}$, forbidden sizes
only $\{2, 3\}$) studied in \cite{SandersGishFourCore}.

%==============================================================
\section{Reproducibility}
\label{sec:repro}
%==============================================================

All six orphans are verified in the foundations module
\texttt{Gen13/targets/foundations/invariants.py}. The full
verification suite (the $48$-invariant module) runs in under $30$
seconds on a standard laptop and reports
all $48$ invariants passing.

%==============================================================
\section*{Acknowledgments}
%==============================================================
The orphans collected here have accumulated over multiple internal
sprints (Sprints 14--18) and through the lens-taxonomy session of
May 2026. We thank the collaborators of those sprints --- in
particular M.~Gish, C.~A.~Luther, B.~Mayes, and H.~J.~Johnson --- for
the structural verifications that produced each orphan. The
foundations module's $48$-invariant verification harness was
developed during the late-April 2026 session and provides the
green-light gate for the bundle.

%==============================================================
\begin{thebibliography}{99}

\bibitem{SandersForcing}
B.~R.~Sanders, M.~Gish,
\emph{The CL Forcing Axioms: A1--A9 Uniquely Force the Canonical Composition Lattice},
Submitted to \emph{Algebraic Combinatorics} (2026)
[J33 of the J-series].

\bibitem{SandersGishThreeSubstrate}
B.~R.~Sanders, M.~Gish,
\emph{The Three-Substrate Architecture: $\mathrm{CL}_\TSML$, $\mathrm{CL}_\BHML$, $\mathrm{CL}_\mathrm{STD}$ as Parallel Substrates},
Submitted to \emph{Algebra Universalis} (2026)
[J31 of the J-series].

\bibitem{SandersGishHARMONYLadder}
B.~R.~Sanders, M.~Gish,
\emph{The 70/71/72/73 HARMONY Ladder: Four Independent Algebraic Constructions},
Submitted to \emph{J.\ Combinatorial Theory A} (2026)
[J30 of the J-series].

\bibitem{SandersGishWobble}
B.~R.~Sanders, M.~Gish,
\emph{Wobble Localization: Prime 11 in TSML\_RAW Characteristic Polynomial},
In preparation [J43 of the J-series].

\bibitem{SandersGishFpBHMLExtension}
B.~R.~Sanders, M.~Gish,
\emph{F_p Extensions of CL\_BHML: Universality Across Six Prime Fields},
Submitted to \emph{Communications in Algebra} (2026)
[J34 of the J-series].

\bibitem{SandersGishFourCore}
B.~R.~Sanders, M.~Gish,
\emph{The 4-Core $\{0,7,8,9\}$: Joint TSML+BHML Closure and the Universal Attractor},
In preparation [target Algebraic Combinatorics, 2026].

\bibitem{SandersTIG}
B.~R.~Sanders,
\emph{The Crossing Lemma at $50\,$Hz: A Substrate-First Theory of Information Generation},
Manuscript in preparation, 2026.

\bibitem{LensCatalog}
B.~R.~Sanders, M.~Gish,
\emph{Variant Catalog of $\Z/10\Z$ Composition Lattices},
Reference compilation, 2026
[\texttt{Atlas/LENS\_TAXONOMY\_2026-05-06/VARIANT\_CATALOG.md}].

\end{thebibliography}

\end{document}
